\origpage{258}{ROBERTSON AND SEYMOUR}
\noindent ends, both with valency at least three. (The \textit{valency} of a vertex is the number of edges incident with it, counting loops twice.) Let $G_2$ be the graph obtained from $G_1$ by contracting $e$. Then we say that $G_1$ is obtained from $G_2$ by \textit{splitting a vertex}. We say that $G_1$ is obtained from $G_2$ by \textit{splitting vertices} if there is a sequence
\begin{equation*}
G_2 = H_0, H_1, H_2, \dots, H_k = G_1
\end{equation*}
of graphs where for $1 \le i \le k$, $H_i$ is obtained from $H_{i-1}$ by splitting a vertex. The lemma we need is the following.

\stmt
(2.2) \textit{If $H$ is a minor of $G$, there is a graph $H'$ which can be obtained from $H$ by splitting vertices and which is topologically contained in $G$.}

\textit{Proof.} Let $G'$ be a minimal subgraph of $G$ which has $H$ as a minor. Then evidently $H$ is obtained from $G'$ by contracting some edges of $G'$, say those in $E_1 \subseteq E(G')$. Now if $e \in E_1$, the graph obtained from $G'$ by deleting $e$ and contracting all other edges of $E_1$ does not have $H$ as a minor, while the one obtained by contracting all edges of $E_1$ does. It follows that $e$ is not a loop of the graph obtained from $G'$ by contracting all edges of $E_1 - \{e\}$. Hence, $E_1$ is the edge set of a subforest of $G'$. For a similar reason, no edge in $E_1$ is incident with a vertex with valency one in $G'$.

Choose a sequence $e_1, \dots, e_n$ of distinct edges of $E_1$ with $n$ maximum, such that for $1 \le i \le n$ there is a vertex of $G'$ of valency two incident both with $e_i$ and with some edge distinct from $e_1, \dots, e_i$. Let $H'$ be the graph obtained from $G'$ by contracting $e_1, \dots, e_n$. We claim that $H'$ satisfies the theorem. Certainly (from the definition of $e_1, \dots, e_n$) $H'$ is topologically contained in $G'$ and hence in $G$; it remains to show that $H'$ can be obtained from $H$ by splitting vertices.

Put $E_2 = E(H') \cap E_1$. Then $E_2$ is the edge-set of a subforest of $H'$, and $H$ may be obtained from $H'$ by contracting the edges in $E_2$. Now no edge $e \in E_2$ is incident with a vertex of $H'$ with valency 2, because we could then set $e_{n+1} = e$, contrary to the maximality of $n$. Thus every edge of $E_2$ has both ends with valency $\ge 3$ in $H'$. Let $E_2 = \{f_1, \dots, f_k\}$, and for $0 \le i \le k$ let $F_i$ be $\{f_j \colon i < j \le k\}$. Let $H_i$ be the graph obtained from $H'$ by contracting all edges in $F_i$. Then $H_0 = H$ and $H_k = H'$, and for $1 \le i \le k$ $H_{i-1}$ is obtained from $H_i$ by contracting $f_i$. The ends of $f_i$ in $H_i$ are distinct, since $E_2$ is the edge-set of a subforest of $H'$. We claim that both ends of $f_i$ have valency $\ge 3$ in $H_i$. For let $u$ be an end of $f_i$ in $H_i$, and let $v$ be the corresponding end of $f_i$ in $H'$. Let $T$ be the component of the forest $(V(H'), F_i)$ with $v \in V(T)$. Now every edge of $H'$ not in $E(T)$ but incident with a vertex of $T$ is incident with $u$ in $H_i$. If $|V(T)| = 1$ it follows that $u$ has valency $\ge 3$ in $H_i$, since $v$ has valency $\ge 3$ in $H'$. If $|V(T)| \ge 2$, there is a vertex $w$ of $T$, with valency 1 in $T$, different from $v$. Then $w$ has valency $\ge 3$ in $H'$, and
